\section{Introduction}
\label{sec:intro}

Decentralized storage networks (DSNs)~\cite{filecoin,sia,shelby,permacoin,retricoin,arweave,walrus,bft-dsn,swarm} provide reliable, long-term data storage without trusting any single participant.
They are increasingly deployed in production:
60\% of decentralized applications use the IPFS/Filecoin~\cite{filecoin} stack for off-chain storage~\cite{filecoin-web3}, Arweave~\cite{arweave} archives ledger data for Solana, Avalanche, and other chains~\cite{arweave-solana,arweave-blockchains}, and the top DSNs collectively host over an Exabyte of data~\cite{filecoin-size}.
Unlike centralized storage, DSNs rely on permissionless nodes that join and leave freely, and use blockchains for governance and coordination in place of trusted operators.

Ideally, a DSN should achieve three properties: \emph{data availability} despite Byzantine participants, \emph{storage scalability} that grows linearly with participants, and \emph{blockchain efficiency} with small, constant on-chain state per data block.
Achieving all three properties, however, is challenging due to Byzantine behaviors.
As in traditional distributed storage, each data block in a DSN is stored redundantly on a \emph{placement group} of nodes.
However, irrational Byzantine nodes can ignore incentives~\cite{bar}, collude within a placement group, and discard stored data simultaneously~\cite{ipfs-attack,bittorrent-attack}.
Data availability therefore depends on sufficient \emph{honest} nodes in each group, yet it is well-known that Byzantine nodes cannot be reliably identified.
A second challenge is erasure coding integrity.
Assigning unique encoded fragments to nodes normally requires coordination, but in a decentralized setting, such coordination demands costly BFT consensus, and a single Byzantine node can poison data read and repair.
We review prior DSN designs in \cref{sec:motivation:review} and conclude that they all fall short of at least one of the three properties (\cref{tab:dsn}).

We trace the root cause to the \emph{data-centric} approach in prior systems.
Existing DSNs form and incentivize each placement group independently, tying rewards to individual data blocks.
To tolerate $\frac{N}{3}$ Byzantine nodes, each group needs more than $\frac{N}{3}$ members, which contradicts storage scalability.
In practice, groups are kept small, allowing Byzantine nodes to concentrate and compromise selected groups.
When they do, honest nodes in those groups receive diluted rewards and may leave, further weakening availability.

Our key insight is that provisioning every group against worst-case Byzantine power is sufficient but not necessary.
We propose a novel \emph{node-centric} approach: instead of allocating incentives per data block, we reward each node equally for its storage contributions, regardless of which groups it joins.
Under this model, Byzantine nodes clustering in a group have no impact on the rewards of honest nodes.
Honest nodes therefore remain in place, and availability is preserved as long as enough of them join each group.

We then design \sys, a concrete DSN instance of this node-centric approach that guarantees sufficient honest nodes in each group.
\sys forms placement group through \emph{random sampling}.
Each node evaluates a verifiable random function (VRF)~\cite{vrf,vrf-rfc} on each data block's hash; the node is \emph{endorsed} for that block's group if the VRF output falls below a public sampling rate.
Because VRF outputs are pseudorandom and independent across nodes, the expected number of honest endorsees per group is uniform and tunable, regardless of how Byzantine nodes distribute and adapt.
We further apply random sampling to erasure coding: each endorsed node independently samples a fragment in the encoding space, without coordination or heavy on-chain metadata.
\sys applies verifiable rateless erasure code~\cite{verify-rateless-ec} to improve space efficiency and to allow any party to check fragment integrity, preventing Byzantine nodes from poisoning data.

We build a complete \sys prototype and prove that it provides strong data availability even in the presence of $\frac{1}{3}$ strong adversaries.
Our \sys design also applies several optimizations that improve the efficiency and cost of placement group discovery and data repair.
%with a target of $N_e = 80$ endorsed honest nodes per group and recovery threshold $k = 32$, the mean time to data loss exceeds 60 years (\cref{sec:design:proof}).
We evaluate the prototype through large-scale simulation and a geo-distributed deployment of 10,000 nodes across five continents.
Overall, this paper makes the following contributions:
\begin{itemize}
    \item We identify the data-centric incentive model as the root cause of the availability-scalability-efficiency trade-off in existing DSNs, and propose a novel node-centric approach that decouples honest-node rewards from Byzantine distribution (\cref{sec:approach}).
    \item We design a VRF-based sampling mechanism for both placement-group formation and erasure coding assignment.
    This mechanism requires only constant on-chain state per data block, supports trustless fragment verification, and tolerates a Byzantine fraction that scales with the network (\cref{sec:design}).
    \item We evaluate \sys extensively: in simulation, \sys maintains data availability over 10 years at $2\times$ redundancy; in real deployment, \sys achieves 30--40 second latency for 1~GB objects, comparable to Swarm~\cite{swarm}, with near-constant performance as the network scales (\cref{sec:eval}).
\end{itemize}
